\documentclass[a4paper,12pt]{article}
\usepackage{amsmath}
\usepackage{graphicx}
\usepackage{multicol}
\usepackage{geometry}
\usepackage{fancyhdr}
\usepackage{tikz}
\usepackage{eso-pic}

\geometry{top=1.5cm, bottom=1.5cm, left=1.5cm, right=1.5cm}

\pagestyle{fancy}
\fancyhf{}
\rhead{Esc. 4-117 Ejército de los Andes}
\lhead{Termodinámica y Máquinas Térmicas}
\cfoot{\thepage}

\AddToShipoutPicture*{
    \put(150,550){\rotatebox{45}{\textcolor[gray]{0.95}{EXAMEN ORIGINAL NO DIGITALIZAR}}}
}

\begin{document}

\begin{center}
    \Large \textbf{Evaluación de Termodinámica y Máquinas Térmicas} \\
    \normalsize Duración: 60 minutos \\
    Alumno/a: \underline{\hspace{6cm}} \hspace{0.5cm} Fecha: \underline{\hspace{2cm}}
\end{center}

\vspace{0.3cm}

\begin{multicols}{2}

\section*{Problemas}

\begin{enumerate}
\item (\textbf{Ley de Charles})\\
Un gas ocupa $2.1\,L$ a $290\,K$. Calcular su volumen final si la temperatura aumenta a $360\,K$, manteniendo constante la presión.

\item (\textbf{Ley de Gay-Lussac})\\
Un gas está a $0.9\,atm$ y $15^\circ C$. Si se lo calienta hasta $95^\circ C$, ¿cuál será su nueva presión si el volumen se mantiene constante?

\item (\textbf{Ecuación general del estado})\\
Un gas ocupa $5.0\,L$ a $1.3\,atm$ y $280\,K$. ¿Qué volumen ocupará a $1.8\,atm$ y $330\,K$?

\item (\textbf{Gas ideal – cálculo de moles})\\
Un recipiente contiene $6.0\,L$ de gas a $2.0\,atm$ y $300\,K$. ¿Cuántos moles de gas hay en el recipiente? ($R = 0.082\,\frac{L \cdot atm}{mol \cdot K}$)

\item (\textbf{Ley de Dalton – presiones parciales})\\
Se mezclan $2\,mol$ de CO\textsubscript{2}, $1\,mol$ de CH\textsubscript{4} y $1\,mol$ de O\textsubscript{2}. Si la presión total es $2.4\,atm$, calcular la presión parcial de cada componente.

\item (\textbf{Ley de Amagat – volúmenes parciales})\\
Tres gases se encuentran en un recipiente: gas A ($3.0\,L$), gas B ($1.0\,L$) y gas C ($2.0\,L$), todos medidos a las mismas condiciones de presión y temperatura. Determinar el volumen total y qué porcentaje ocupa cada gas en la mezcla.

\end{enumerate}

\vspace{0.2cm}

\end{multicols}

\end{document}
